\section{多重比较改变错误的计量单位}
\label{sec:11-Mathematical-Statistics/05-Hypothesis-Testing/04}

实验看板上有十四个指标：一个主指标，五个次级指标，八个护栏指标。这次改动其实什么也没改变——假设我们上帝视角地知道这一点。看板上出现至少一个 \(p<0.05\) 的绿色或红色标记的概率是 \(1-0.95^{14}\approx 0.51\)。一半的场次会有个指标跳出来，然后有人围绕它讲一个故事。

每个 \(p\) 值都算对了，每个检验的水平都是 \(0.05\)。出问题的是承诺的对象：\(\alpha\) 承诺的是``这一个假设被误报的概率''，而看板上要的是``这块板子上有没有一处误报''。计量单位从每次检验换成了整族检验，原来的承诺不再覆盖新的问题。

\subsection{做二十次检验，至少错一次的概率是 \texorpdfstring{\(0.64\)}{0.64}}

\(m\) 个相互独立、且全部为真的原假设，各自用水平 \(\alpha\) 检验，至少出现一个假阳性的概率是

\[
1-(1-\alpha)^{m}.
\]

\(\alpha=0.05\) 时，\(m=1\) 给 \(0.05\)，\(m=14\) 给 \(0.51\)，\(m=20\) 给 \(0.64\)，\(m=100\) 给 \(0.994\)。

\begin{center}
\begin{tikzpicture}[x=0.185cm,y=3.1cm,font=\footnotesize,>=stealth]
  \draw[black!70] (0,0) -- (42,0);
  \draw[black!70] (0,0) -- (0,1.06);
  \foreach \y in {0,0.25,0.5,0.75,1.0}
    { \draw[black!60] (-0.6,\y) -- (0,\y); \node[anchor=east,black!70] at (-1.0,\y) {\y}; }
  \foreach \x in {1,10,20,30,40}
    { \draw[black!60] (\x,-0.02) -- (\x,0); \node[anchor=north,black!70] at (\x,-0.03) {\x}; }
  \draw[black!30,densely dotted] (0,0.642) -- (20,0.642);
  \draw[black!30,densely dotted] (20,0) -- (20,0.642);
  \draw[black!80,thick] plot[domain=0:41,samples=90] (\x,{1-exp(-0.0512933*\x)});
  \draw[black!75,densely dashed] (0,0.0488) -- (41,0.0488);
  \fill[black!85] (20,0.642) circle (1.6pt);
  \node[anchor=west,black!85] at (21,0.55) {\(m=20\)：\(0.64\)};
  \node[anchor=west,black!85] at (30,0.70) {每个用 \(\alpha\)};
  \node[anchor=west,black!75] at (26,0.13) {每个用 \(\alpha/m\)};
  \node[anchor=north,black!70] at (21,-0.16) {检验个数 \(m\)};
  \node[anchor=south,black!70,rotate=90] at (-4.2,0.35) {至少一个假阳性的概率};
\end{tikzpicture}

\emph{检验个数一多，``至少错一次''的概率迅速逼近 \(1\)；把每个检验的水平压到 \(\alpha/m\)，这条曲线被摁回原来的高度。}
\end{center}

图里那条几乎贴着底部的虚线，是把每次检验的水平改成 \(\alpha/m\) 之后的同一个量。它说明补救是可能的，也说明补救不是白来的：每一次检验的判据都被拧紧了。

族错误率（family-wise error rate）把这件事写成定义。设 \(V\) 是被错误拒绝的真原假设个数，

\[
\mathrm{FWER}=P(V\ge 1),
\]

Bonferroni 方法给每个检验分配水平 \(\alpha/m\)，理由是一行 union bound：

\[
P\Bigl(\bigcup_{j=1}^{m}A_j\Bigr)\le\sum_{j=1}^{m}P(A_j)\le m\cdot\frac{\alpha}{m}=\alpha ,
\]

其中 \(A_j\) 是第 \(j\) 个真原假设被误拒的事件。条件对账：这里只用了概率的次可加性，不需要检验之间独立，也不需要知道依赖结构；真原假设的个数未知时用总数 \(m\) 代替，只会更保守。保守的来源也在同一处——检验之间高度相关时（同一批用户的十几个指标通常如此），并集远小于各项之和，Bonferroni 把一笔从未发生的账也付了。

Holm 方法把 \(p\) 值从小到大排，第 \(k\) 个与 \(\alpha/(m-k+1)\) 比较，一旦某个不通过就停止。它在同样宽松的条件下控制 FWER，拒绝集合总包含 Bonferroni 的，功效不会更差，没有理由不用它。

功效的账要单独算一遍。二十个指标做 Bonferroni 校正，单个检验的水平从 \(0.05\) 降到 \(0.0025\)，双侧临界值从 \(1.96\) 抬到 \(3.02\)。按上一节末尾的样本量公式，要维持同样的功效，样本量需要乘上 \(\bigl((3.02+0.84)/(1.96+0.84)\bigr)^{2}\approx 1.9\)。看板上多挂十几个指标，实验就要多跑将近一倍的量——这笔钱要么事前花在样本量上，要么事后花在假阳性上，没有第三种付法。

\subsection{大规模筛选换一种计量方式}

十几个指标的看板可以承受 Bonferroni。一万列特征的筛选不行：\(\alpha/10^{4}=5\times 10^{-6}\)，除了效应极强的少数几列，其余全部被压住，真信号也一起没了。这类场合下，``一个假阳性都不许有''本来也不是合理的诉求——筛选的产出是一份候选清单，清单里混进几个错的可以接受，只要比例可控。

于是把计量单位再换一次。设 \(R\) 是拒绝总数，\(V\) 是其中的假阳性数，错误发现率定义为

\[
\mathrm{FDR}=\E\!\left[\frac{V}{\max(R,1)}\right].
\]

Benjamini--Hochberg 过程把 \(p\) 值排序为 \(p_{(1)}\le\cdots\le p_{(m)}\)，找出满足 \(p_{(k)}\le \frac{k}{m}q\) 的最大下标 \(k^\star\)，拒绝前 \(k^\star\) 个。在独立或某类正依赖条件下，它控制 FDR 不超过 \(q\)。

阈值 \(\frac{k}{m}q\) 随排名放宽这一点，正好对应两种计量方式的差别。Bonferroni 给每个检验固定的额度 \(\alpha/m\)，不管最后发现了多少；BH 的额度随发现数增长——发现得越多，能容忍的假阳性绝对数越多，因为受控的是比例。一万列里若有二百列真有信号，\(q=0.1\) 的 BH 通常能捞回其中相当一部分，而交出的清单里约有一成是错的。清单要拿去做下一轮验证时，这个交易很划算；清单要直接上线时，就不划算。

选择的依据是决策的形状。任何一个假阳性都代价高昂——药物上市、风控策略全量、对外发布的结论——用 FWER；产出是待验证的候选集、后面还有一道独立的关卡——特征筛选、基因组扫描、A/B 实验的探索性指标——用 FDR。

顺带说一句常被忽略的方向问题。护栏指标希望的是``不显著''，即不要发现损害。对护栏指标做多重校正，等于抬高发现损害的门槛，降低了发现损害的功效，与设置护栏的初衷相反。护栏应当反过来处理：用等价性检验，或者干脆放松水平换取功效。

\subsection{家族由决策边界定义，不由表格行数定义}

哪些检验算在同一族里，取决于哪些结果会共同支撑一个决定，以及最终会被一起看到多少。做了五十个分析只把最小的三个写进报告，家族是五十，不是三。这条规则没有形式化的判据，但有个可操作的替身：把``从数据到这份报告''这条路径上所有被数据影响过的选择数出来。

多重性的来源比指标个数广得多。多个子群（新用户、老用户、各渠道、各机型），多个时间窗（首日、七日、三十日），多种模型设定与协变量组合，多种离群点剔除规则，事后挑选的单侧方向，边跑边看的中途查看（上一节那个可选停时的漏洞），先用同一份数据筛特征再用它做检验，以及多个实验只报成功的那个。这些选择每一个都在暗中扩大家族，而报告表格上一行也看不出来。名义 \(p\) 值和名义区间的校准，在这些操作之后都不再成立。

区间也受同样的影响。想让一组区间同时以 \(1-\alpha\) 的概率覆盖全部参数，最简单的办法仍是 Bonferroni，把每个区间的置信水平提到 \(1-\alpha/m\)，区间随之变宽；线性模型的全部对比、时间序列的置信带，可以利用相关结构做得比这更紧。更要紧的是另一件事：只给显著的那些结果报普通区间，区间本身已经带上了选择偏差——被选中的条件是估计值足够大，所以条件分布的中心比真值偏高。这与低功效实验的夸大比是同一个机制，在筛选场景里更严重。需要无偏的数字时，办法是条件化在选择事件上做推断，或者更省事地，留一份独立数据专门用来确认。

\subsection{校正是补救，分工才是解法}

事后校正处理的是已经发生的多重性，它管不了那些没有被记录下来的选择。若假设是在同一份数据上反复试探生成的，把最后剩下的检验个数代进 Bonferroni，覆盖不了完整的搜索过程——分母根本不知道自己应该是多少。

可靠的做法是在设计阶段就把探索与确认分开。预先写定主指标和假设家族，主指标只有一个，次级指标列在文档里且说明用途；探索阶段放开跑，用 FDR 控制清单质量，明确标注结论是待验证的；确认阶段用独立的数据或新一轮实验，只检验从探索阶段带出来的那几条，家族小，可以用 FWER。所有分析路径都记下来，包括没有出结果的那些，否则后来的人无法判断家族有多大。

回到开头那块看板。改法不复杂：主指标一个，判定上不上线只看它；次级指标用 Holm 校正后一起报，写清它们不参与决策；护栏指标改成等价性检验，声明容忍界；实验跑完之前不看结论。这些约定都写在开跑之前，事后不改。做到这一步，看板上那半数场次的假信号就不会再变成任何人的季度目标。

这一单元从一条决策规则的两类错误出发，走到多重比较下错误率的重新计价，中间反复出现同一个要求：分析流程必须先于数据固定。当流程本身也需要用数据来挑选——选模型、选超参数、选特征——就得有一套办法把``评估''和``选择''隔开。这是下一单元\hyperref[ch:091]{``重采样与模型选择：让数据评估整个流程''}要处理的问题。
